\chapter{根文件系统与 devfs}
\label{chap:rootfs}
本章对应当前 miniOS 的实现，说明内核如何从 SD 卡找到并挂载根文件系统（xv6 血统的 minifs）、VFS 挂载表的形状、\file{/boot} 的 FAT32 交接（bootfs）、\file{/dev} 设备文件（devfs），以及用户态 \file{/init}、\file{/bin/sh}是如何从根文件系统里跑起来的。

\begin{note}
\textbf{更新}：\file{/boot} 已改为内核直接读写的 FAT32 文件系统（bio + fat32），挂载顺序也改成先 \file{/} 再 \file{/boot}，详见 第~\ref{chap:bootfat32} 章。本章的 bootfs 小节现在描述的是 FAT32 挂载失败时的后备路径。
\end{note}

当前已打通的链路：

\begin{itemize}
  \item SPL 读 SD 128 KiB 的 bootpack 装内核，再读 FAT32 分区把 XR829固件装进 DDR 并补写 bootpack 头（交接）；
  \item 内核 \code{kernel\_boot()} 依次挂载 \file{/}（minifs，MBR 第 2 分区按需读）、\file{/boot}（FAT32，失败时退回 bootfs）、\file{/dev}（devfs），任一最终失败即 PANIC；
  \item XR829 驱动从 \file{/boot/*.bin} 经 VFS 读固件（不再打进 bootpack）；
  \item 内核从 \file{/} 读出 \file{/init} 创建第一个用户进程；\code{init} fork+exec \file{/bin/sh}，shell 及 \code{ls/cat/echo/...} 都从 minifs 经 VFS 加载。
\end{itemize}

尚未实现：文件写入（minifs 只读）、minifs 日志区、devfs 作为用户态stdio 后端、\code{mount}/\code{umount} 系统调用、多个块设备。

\section{SD 卡布局}

\begin{table}[H]
\centering
\small
\begin{tabularx}{\linewidth}{l>{\raggedright\arraybackslash}X>{\raggedright\arraybackslash}X}
\toprule
区域 & 位置 & 内容 \\
\midrule
Boot0（SPL） & 8 KiB 偏移，512 B & miniBoot，SRAM 里跑，负责 DDR init + 装内核 \\
bootpack & 128 KiB 偏移 & 512 B header + kernel 镜像（见 \code{bootpack.h}） \\
MBR 分区 1 & FAT32 & \code{boot\_xr829.bin} / \code{fw\_xr829.bin} / \code{sdd\_xr829.bin}（macOS 挂载于 \file{/Volumes/bootfs}） \\
MBR 分区 2 & minifs 根文件系统 & \code{build-t113/minifs.img} 整镜像 dd 进去，≥ 32 MiB \\
\bottomrule
\end{tabularx}
\end{table}

分区 2 没有文件系统以外的含义：\code{rootdev.c} 只是取 MBR 第二个分区表项的 LBA/长度当块设备用。分区本身用任何分区工具建（类型 0x83 即可），FAT32 分区 1 是给 SPL 固件加载和人工换固件用的。

\begin{txtcode}
SD 8 KiB    +------------+
            | Boot0 SPL  |  读 128 KiB bootpack → 装内核
SD 128 KiB  +------------+
            | bootpack   |  header + kernel @0x40200000
分区 1      +------------+
            | FAT32      |  SPL 读固件 → DDR 0x40100000/0x40110000/0x40140000
分区 2      +------------+
            | minifs     |  内核 rootdev 按需读（块设备，非预载）
            +------------+
\end{txtcode}

\section{总调用链}

\begin{txtcode}[basicstyle=\ttfamily\scriptsize]
miniBoot SPL
+-- spl_main()
    +-- load_bootpack()
        +-- t113_sd_init()                    SMHC0 识别卡
        +-- 读 SD 128 KiB bootpack header     CRC32 校验
        +-- 逐条目读镜像到 load_addr + CRC32  kernel → 0x40200000
        +-- header 拷到 0x401ff000            内核侧查找入口
        +-- load_fat32_firmware()             FAT32 交接
            +-- fat32_mount()                 MBR 找 FAT32 分区, BPB 解析
            +-- fat32_load_root() × 3         boot/fw/sdd_xr829.bin → DDR
            +-- bootpack_slot()               补 3 个条目(offset=0 常驻), 重算 header CRC
        +-- 跳 0x40200000

miniOS kernel_main()                          boards/100ask_t113/platform_kernel.c
+-- t113_uart_init() / t113_mmu_init()
+-- t113_bootpack_dump()                      打印全部条目（含 SPL 交接的固件）
+-- kernel_boot()                             kernel/init/main.c
    +-- t113_page_alloc_init()                页分配器
    +-- t113_vfs_init()                       mount/super_block/inode/dentry/file 池清零
    +-- t113_minifs_mount_rootdev()           → /       失败 PANIC
    |   +-- t113_rootdev_init()               MBR 分区 2 解析
    |   |   +-- t113_sd_init()                SMHC0, 1-bit, 轮询
    |   |   +-- 读 LBA0, 验 0x55AA, 取表项 2  type/lba/sectors
    |   +-- 读 root_lba+2 扇区                块 1 = superblock
    |   +-- 验 magic 0x10203040, size×2 ≤ 分区扇区数 → t113_vfs_mount_fs("/")
    +-- t113_fat32_mount_vfs()                → /boot   读写 FAT32
    |   +-- 失败 → t113_bootfs_mount_vfs()    → /boot   SPL 交接视图，再失败 PANIC
    +-- t113_devfs_init()                     → /dev    失败 PANIC
    +-- net_init()                            CONFIG_NET
    +-- t113_driver_core_init()               平台/sdio 驱动绑定
    |   +-- uart probe → t113_chrdev_register(1/3, 0) + t113_devfs_create_chrdev("console"/"ttyS0")
    |   +-- xradio probe → need_image("/boot/*.bin")  经 VFS 读 FAT32 上的固件
    +-- t113_sched_init()
    +-- start_init_process()
    |   +-- t113_vfs_open("/init") → t113_vfs_read() → t113_create_init_user_process()
    +-- Timer0 + 开中断                       之后 idle: wfi

用户态
+-- /init (init.c)
    +-- fork() + exec("/bin/sh")              sys_exec: VFS 打开 → ELF → 新地址空间
    +-- wait() 孤儿回收, shell 退出后重启
\end{txtcode}

\file{/} 和 \file{/dev} 是硬依赖；\file{/boot} 先试 FAT32、再退回 bootfs，两者都失败才算失败。任何一个最终失败，\code{kernel\_boot()} 都打印 PANIC 并 \code{wfi} 死循环——没有根文件系统就没有用户态，无处可退。这对应 Linux “unable to mount root fs” 的 panic，只是 miniOS 连 initramfs 兜底都没有，也不需要：内核本身在 bootpack 里，卡上就是它的全部世界。挂载过程的图解见下一节。

\section{图解：内核怎样把根文件系统挂到 \file{/}}

\subsection{挂载结果}

先纠正一个容易混的点：devfs 挂在 \file{/dev}，不是 \file{/boot}。\file{/boot} 挂的是 SD 卡第一个分区上的 FAT32；只有 FAT32 挂载失败时，才退回到 bootfs（SPL 预加载镜像的只读视图）。启动完成后，整个命名空间是这样的：

\begin{txtcode}[basicstyle=\ttfamily\scriptsize]
SD 卡                                  VFS 挂载表 mounts[8]          命名空间
+--------------+
| MBR          |                       [0] "/"     → minifs_vfs_ops   /
| Boot0 / pack |                       [1] "/boot" → fat32_ops        +-- init
+--------------+  第 1 个 FAT32 分区     (失败时  → bootfs_ops)        +-- bin/ sh ls cat ...
| 分区 1 FAT32 | --------------------------------------------------→  +-- boot/                FAT32
+--------------+  MBR 表项 2                                          |   +-- boot_xr829.bin
| 分区 2 minifs| --------------------------------------------------→  |   +-- fw_xr829.bin
+--------------+                       [2] "/dev"  → devfs_ops        |   +-- sdd_xr829.bin
                                                                      |   +-- wifi.txt
内存：devfs nodes[16] -----------------------------------------------→ +-- dev/                 devfs
       （uart 驱动 probe 时注册）                                          +-- console
                                                                          +-- ttyS0
\end{txtcode}

\code{find\_mount()} 按最长前缀选后端。路径分隔符要求完整匹配，所以 \file{/bootx} 不会落到 \file{/boot}：

\begin{table}[H]
\centering
\small
\begin{tabularx}{\linewidth}{>{\raggedright\arraybackslash}X>{\raggedright\arraybackslash}X>{\raggedright\arraybackslash}X}
\toprule
路径 & 匹配的挂载点 & 后端 \\
\midrule
\file{/init} & \file{/} & minifs \\
\file{/bin/sh} & \file{/} & minifs \\
\file{/boot} & \file{/}、\file{/boot}，取 \file{/boot} & fat32（根目录） \\
\file{/boot/fw\_xr829.bin} & \file{/boot} & fat32 \\
\file{/bootx} & \file{/} & minifs（找不到就返回失败） \\
\file{/dev/console} & \file{/dev} & devfs \\
\bottomrule
\end{tabularx}
\end{table}

\subsection{启动时的挂载顺序}

\begin{txtcode}[basicstyle=\ttfamily\scriptsize]
SPL 跳到 0x40200000
+-- kernel_entry.S
    +-- kernel_main(dram_size)                     boards/100ask_t113/platform_kernel.c
        +-- t113_uart_init()                       早期串口
        +-- t113_mmu_init()                        内核 L1 表，打开 MMU
        +-- t113_bootpack_dump()
        +-- kernel_boot(pool_base, dram_end)        kernel/init/main.c，与板子无关的部分
            +-- t113_page_alloc_init()
            +-- t113_vfs_init()                    清空 mount/super_block/inode/dentry/file 五个池
            +-- t113_minifs_mount_rootdev()        (1) "/"   失败 → PANIC
            |   +-- t113_rootdev_init()
            |   |   +-- t113_sd_init()             SMHC0：CMD0/8/ACMD41/2/3/7，切到 12 MHz
            |   |   +-- t113_sd_read(0, 1, mbr)    验 55 AA
            |   |   +-- 表项 2：root_lba、root_sectors_count
            |   +-- t113_rootdev_read(2, 2, block_cache)     minifs 块 1 = 超级块
            |   +-- 验 magic 0x10203040，size × 2 ≤ 分区扇区数
            |   +-- t113_vfs_mount_fs("/", MINIFS_MAGIC, &fs_sb, &minifs_vfs_ops)
            +-- t113_fat32_mount_vfs()             (2) "/boot"
            |   +-- binit()                        建 bcache
            |   +-- bread(0)：MBR，找第一个 FAT32 类型的分区
            |   +-- bread(part_lba)：BPB 校验并算出布局
            |   +-- t113_vfs_mount_fs("/boot", 0x4d44, &fs, &fat32_ops)
            |       （先 lookup_path("/boot") 取 minifs 里的挂载点目录）
            |   失败 → t113_bootfs_mount_vfs()     (2') 退回 bootfs，再失败 → PANIC
            +-- t113_devfs_init()                  (3) "/dev"  失败 → PANIC
            +-- net_init()                         CONFIG_NET
            +-- t113_driver_core_init()            设备、驱动绑定
            |   +-- uart probe → t113_chrdev_register(1, 0, &console_fops)
            |   |                t113_chrdev_register(3, 0, &console_fops)
            |   |                t113_devfs_create_chrdev("console", 1, 0)
            |   |                t113_devfs_create_chrdev("ttyS0", 3, 0)
            |   +-- mmc_rescan_all() → xradio probe → need_image("/boot/*.bin")
            +-- t113_sched_init()
            +-- start_init_process()
            |   +-- t113_vfs_open("/init") → minifs
            |   +-- t113_vfs_read() 整个 ELF 读进临时页
            |   +-- t113_create_init_user_process()
            +-- 打开定时器和中断，进入 wfi 循环
\end{txtcode}

先挂 \file{/} 有一个具体原因：\code{t113\_rootdev\_init()} 负责 \code{t113\_sd\_init()}，也就是初始化 SMHC0 和 SD 卡；FAT32 和它共用这个控制器。驱动必须在\file{/boot}、\file{/dev} 都挂好之后才初始化，因为 XR829 的 probe 要从 \file{/boot} 读固件，UART 的 probe 要往 \file{/dev} 里注册节点。

\subsection{rootdev：分区到块设备}

minifs 的块是 1 KiB，SD 扇区是 512 B，\code{rootdev.c} 只做一次线性换算：

\begin{txtcode}
minifs 块号 b  --→  分区内扇区 2b, 2b+1  --→  整卡 LBA root_lba + 2b

t113_rootdev_read(sector, count, dst)
+-- 越界检查：sector + count ≤ root_sectors_count
+-- t113_sd_read(root_lba + sector, count, dst)
\end{txtcode}

它相当于 Linux 的分区块设备（\file{/dev/mmcblk0p2}），或 xv6 的 \code{ROOTDEV}加上 \code{virtio\_disk\_rw()}。区别是它不经过任何缓存，缓存放在 minifs 里（只缓存一块）。

\subsection{minifs 的磁盘布局和查找}

minifs 沿用 xv6 \code{mkfs} 的格式（\code{MINIFS\_MAGIC} 就是 xv6 的 \code{FSMAGIC} 0x10203040），布局和 xv6 一样：

\begin{txtcode}[basicstyle=\ttfamily\scriptsize]
块:  0       1        2 ... logstart+nlog-1   inodestart ...     bmapstart ...   数据块 ...
    +------+--------+----------------------+----------------+-------------+-------------+
    | boot | super  | log（minifs 不用）    | inode 块        | 空闲位图     | 文件/目录内容 |
    +------+--------+----------------------+----------------+-------------+-------------+
           magic size nblocks ninodes nlog logstart inodestart bmapstart

struct minifs_dinode（64 B，每块 16 个）
+-- type  1 = 目录，2 = 文件
+-- major minor nlink
+-- size
+-- addrs[12] 直接块 + addrs[12] 一个间接块（256 项）  → 最大 268 KiB

struct minifs_dirent（32 B）
+-- inum（u16）
+-- name[30]
\end{txtcode}

\code{exec("/bin/sh")} 在 minifs 里的查找：

\begin{txtcode}
t113_vfs_open("/bin/sh")
+-- minifs_vfs_lookup()
    +-- t113_minifs_open() → namei("/bin/sh")
        +-- current = ROOTINO (1)
        +-- "bin": directory_lookup(1, "bin")
        |   +-- read_inode(1)          块 inodestart + 1/16，取第 1 项
        |   +-- 逐个 32 B 目录项 inode_read() → 找到 inum = N
        +-- "sh":  directory_lookup(N, "sh") → inum = M
        +-- *inum = M；read_inode(M) 得到 size
sys_exec → t113_vfs_read(file, image, size)
+-- minifs_vfs_read() → t113_minifs_read(M, off, dst, n)
    +-- inode_read()
        +-- 每块：inode_block(inode, off / 1024)       直接块或间接块
                  block_address(b)                     命中 1 块缓存就直接用
                  +-- t113_rootdev_read(2b, 2, block_cache)
\end{txtcode}

和 xv6 的 \code{namex()} 相比，minifs 只读，所以不需要 inode 缓存、\code{ilock}、引用计数，也不需要日志。每次查找都从根目录开始，因为 T113 没有每进程的当前目录 inode（\code{getcwd} 返回的是固定的 \file{/}）。

\subsection{devfs：\file{/dev} 下的字符设备}

\begin{txtcode}[basicstyle=\ttfamily\scriptsize]
驱动 probe 时                           用户或内核 open("/dev/console")
t113_chrdev_register(1, 0,              t113_vfs_open("/dev/console", RDWR)
                     &console_fops)     +-- lookup_path：find_mount → "/dev" 挂载根 dentry
+-- chrdev_table[i] = { dev, ops }      |   +-- 分量 "console" → devfs_lookup("/dev/console")
t113_devfs_create_chrdev(               |       +-- devfs_basename() → "console"
   "console", 1, 0)                     |       +-- 遍历 nodes[]，id = MINIOS_MKDEV(1, 0)
+-- nodes[i] = { name, major, minor }   |           type = CHAR，private_data = &nodes[i]
                                        +-- devfs_open()：只接受 RDONLY / RDWR
                                        t113_vfs_write(f, "hi", 2)
                                        +-- devfs_write() → t113_chrdev_write(0x10000)
                                            +-- chrdev_table 按设备号 → console_fops.write()
                                                → UART3 发送
\end{txtcode}

devfs 的节点不在任何磁盘上，是驱动在内存里登记的表。这和 Linux 的 devtmpfs同一个思路：驱动注册设备时，内核自动在 \file{/dev} 下建节点。xv6 的做法不同：设备节点是磁盘上的 inode（\code{T\_DEVICE}，\code{mknod("console", CONSOLE, 0)} 由\code{init} 创建），\code{read/write} 时按 \code{ip->major} 查 \code{devsw[]} 分派。

用户态 \file{/init} 打开 \file{/dev/console}，再用 \code{dup()} 建立 fd 0/1/2；所有控制台I/O 都走 VFS → devfs → chrdev → UART，不再在系统调用层特判标准 fd。

\subsection{与 xv6、Linux 的挂载方式对照}

\begin{table}[H]
\centering
\small
\begin{tabularx}{\linewidth}{l>{\raggedright\arraybackslash}X>{\raggedright\arraybackslash}X>{\raggedright\arraybackslash}X}
\toprule
方面 & xv6 & Linux & T113 miniOS \\
\midrule
根文件系统从哪来 & 固定 \code{ROOTDEV}（virtio 盘，整盘是一个 fs） & \code{root=} 参数，\code{prepare\_namespace()} → \code{mount\_root()}，常先经 initramfs & MBR 第 2 分区，\code{t113\_minifs\_mount\_rootdev()} \\
何时挂载 & 第一个进程的 \code{forkret()} 里 \code{fsinit(ROOTDEV)}，因为读盘要 sleep & \code{kernel\_init} 线程里，在 exec \file{/sbin/init} 之前 & \code{kernel\_boot()} 里，在创建 init 之前；读盘是轮询，不用 sleep \\
有没有挂载表 & 没有，只有一个文件系统 & \code{struct mount} 树，\code{vfsmount} 挂在 dentry 上 & \code{vfs\_mount} 父子关系、mountpoint dentry，路径选择暂兼容最长前缀 \\
第二个文件系统 & 不支持 & 任意 & \file{/boot} FAT32、\file{/dev} devfs \\
设备文件 & 磁盘上的 \code{T\_DEVICE} inode + \code{devsw[major]} & devtmpfs，驱动注册时建节点 & devfs，驱动注册时登记 \\
挂载失败 & \code{panic("invalid file system")} & \code{VFS: Unable to mount root fs} panic & \code{kernel\_platform\_panic("cannot mount root filesystem")} \\
用户态 mount & 无 & \code{mount(2)} & 无 \\
\bottomrule
\end{tabularx}
\end{table}

\subsection{T113 实现上的区别}

\begin{itemize}
  \item 两个文件系统共用一个 SD 控制器（SMHC0）和一张卡。\code{rootdev.c} 不经\code{bio.c}；\code{fat32.c} 经 \code{bio.c}。两条路径都在 \code{t113\_sd\_read/write} 里的\code{sd\_lock} 汇合。
  \item VFS 已建立 \code{super\_block → dentry → inode → file} 对象链、引用计数和 inode/ dentry 缓存，并让子 mount 保存父 mount 与 mountpoint。minifs/FAT32/devfs仍通过兼容适配层接收累计的绝对路径；下一阶段可逐个改成真正的“父 inode + 单个分量” lookup，之后再完整处理 \code{.}、\code{..} 和跨 mount 回退。
  \item 卡上的 SPL、bootpack、两个分区由同一个 \code{make install} 写入；分区 2 必须足够放下 \code{minifs.img}（≥ 32 MiB）。
\end{itemize}

\section{VFS：super\_block、inode、dentry}

VFS 已改成 Linux 式的对象模型，完整说明见 第~\ref{chap:vfs} 章。这里只列和挂载、根文件系统直接相关的几点：

\begin{table}[H]
\centering
\small
\begin{tabularx}{\linewidth}{>{\raggedright\arraybackslash}X>{\raggedright\arraybackslash}X}
\toprule
miniOS & Linux \\
\midrule
\code{vfs\_mount} + \code{parent/mountpoint/mnt\_root} & \code{struct mount} / \code{vfsmount} \\
\code{super\_block} + \code{super\_operations}，\code{s\_backend} 指向旧后端 & \code{struct super\_block} \\
\code{dentry\_cache[128]}、\code{inode\_cache[128]} 与引用计数 & dcache / icache \\
\code{file} + \code{file\_operations} + \code{f\_pos\_lock} & \code{struct file} + \code{f\_count} \\
全局 \code{vfs\_lock} 一把互斥 & per-inode 锁 + RCU \\
\bottomrule
\end{tabularx}
\end{table}

\begin{itemize}
  \item \textbf{挂载}：\code{t113\_vfs\_mount\_fs(path, magic, fs\_info, ops)} 为每个文件系统建一个\code{super\_block}、一个根 inode 和一个根 dentry；不是 \file{/} 时还要先在父文件系统里 \code{lookup\_path(path)} 找到挂载点目录，所以 minifs 镜像里必须有 \file{/boot}、\file{/dev} 两个空目录（\code{mkfs\_minifs} 已经建好）。
  \item \textbf{选文件系统}：仍按最长前缀匹配（\file{/bootx} 不会落到 \file{/boot}），然后从那个挂载的根 dentry 开始逐级建 dentry。
  \item \textbf{后端接口没变}：\code{vfs\_backend\_ops.lookup} 仍收到累计的绝对路径，由 \code{vfs.c}里的 \code{backend\_iops/backend\_fops} 把 inode/file 操作转成后端调用。
  \item \textbf{身份}：minifs 的 inode 号、FAT32 的 \code{fat\_inode} 槽位地址、bootfs 的 DDR地址、devfs 的设备号成为 \code{inode->i\_ino}；同一文件多次打开共享一个 inode。
  \item 读写带 \code{f\_pos} 语义，但\textbf{没有 lseek}；标志支持 \code{VFS\_O\_RDONLY} / \code{VFS\_O\_RDWR}，外加 FAT32 用的 \code{VFS\_O\_CREATE} / \code{VFS\_O\_TRUNC}（值取自 xv6 \code{fcntl.h}）。\code{write} 由 devfs 和 FAT32 实现；minifs 和 bootfs 只读。
  \item \textbf{注意}：\code{minifs\_vfs\_lookup()} 目前只接受普通文件，逐级查找要它返回目录\file{/boot}、\file{/dev}、\file{/bin}，所以挂载 \file{/boot} 会失败，详见第~\ref{chap:vfs} 章。
\end{itemize}

\section{rootdev：MBR 分区 2 块设备}

\code{src/kernel/fs/rootdev.c} + \code{src/drivers/mmc/sunxi-sd.c}：

\begin{itemize}
  \item \code{sunxi-sd.c} 是 \textbf{SPL 和内核共用的一份源码}（\code{MINIOS\_KERNEL} 宏区分是否加 mutex）：SMHC0（\code{0x04020000}，PF0..PF5，1-bit 总线），CMD17 单块读、轮询、CPU 搬 FIFO。SPL 用它读 bootpack/FAT32，内核用它读根文件系统——同一张卡的同一个控制器。
  \item \code{t113\_rootdev\_init()}：读 LBA 0，验证 \code{0x55 0xAA} 签名，取 MBR偏移 \code{446+16} 的第二个分区表项：type（偏移 4）、起始 LBA（+8）、扇区数（+12），type 非零、LBA 非零、≥ 4 扇区才接受。
  \item \code{t113\_rootdev\_read(sector, count, dst)} 做范围检查后转\code{t113\_sd\_read(root\_lba + sector, ...)}。
\end{itemize}

\section{minifs：xv6 磁盘 ABI 的只读实现}

\code{tools/mkfs\_minifs.c}（主机侧）明确保留 xv6 mkfs 的磁盘格式，内核侧\code{src/kernel/fs/minifs.c} 是只读遍历实现。常量与 xv6 一致：

\begin{table}[H]
\centering
\small
\begin{tabularx}{\linewidth}{ll>{\raggedright\arraybackslash}X}
\toprule
常量 & 值 & 说明 \\
\midrule
\code{MINIFS\_BSIZE} & 1024 & 块大小（= 2 个 SD 扇区） \\
\code{MINIFS\_MAGIC} & 0x10203040 & xv6 \code{FSMAGIC} \\
\code{MINIFS\_ROOTINO} & 1 & 根目录 inode 号 \\
\code{MINIFS\_NDIRECT} & 12 & 直接块 \\
\code{MINIFS\_DIRSIZ} & 30 & 文件名上限 \\
\code{dinode} & 64 B & type/major/minor/nlink/size/addrs[13] \\
\code{dirent} & 32 B & u16 inum + name[30] \\
\bottomrule
\end{tabularx}
\end{table}

镜像布局（\code{FSSIZE=32768} 块 = 32 MiB，\code{NINODES=200}）：

\begin{txtcode}
块 0            空（xv6 惯例留空）
块 1            superblock
块 2..31        log 区（30 块，只读实现不使用，格式保留）
块 32..44       inode 区（200/16+1 = 13 块，每块 16 个 dinode）
块 45..49       位图（32768/8192+1 = 5 块）
块 50..         数据区（含目录内容）
\end{txtcode}

mkfs 产出的目录树（\code{miniOS/Makefile} 打包 \code{\_init \_sh \_ls \_echo \_pwd \_cat \_sleep \_uptime}）：

\begin{txtcode}
/                ino 1
+-- bin/         sh, ls, echo, pwd, cat, sleep, uptime
+-- dev/         空目录（minifs 里的占位；真设备在 devfs）
+-- init         与 /bin/init 同一 inode（nlink=2，硬链接）
\end{txtcode}

单文件上限 = 12 直接块 + 256 间接块 = 268 块 = 268 KiB；xv6 风格目录 inode 的 \code{size} 被取整到完整块（\code{round\_directory()}）。

\subsection{两种挂载方式}

\begin{table}[H]
\centering
\small
\begin{tabularx}{\linewidth}{>{\raggedright\arraybackslash}X>{\raggedright\arraybackslash}X>{\raggedright\arraybackslash}X}
\toprule
函数 & 数据来源 & 用途 \\
\midrule
\code{t113\_minifs\_mount\_vfs(image, bytes)} & DDR 里的镜像（\code{minifs\_image.o} 可 objcopy 进内核） & 备用路径；T113 当前未链接 \\
\code{t113\_minifs\_mount\_rootdev()} & MBR 分区 2，按需读块 & \textbf{当前使用} \\
\bottomrule
\end{tabularx}
\end{table}

块设备模式的关键是\textbf{单块缓存}：\code{block\_cache[1024]} + \code{cached\_block}只缓存最近一块。\code{block\_address(block)} 命中缓存直接返回指针，未命中则 \code{t113\_rootdev\_read(block*2, 2, cache)} 读两个扇区。这等于一个只有一项的 page cache——读 \file{/init} 这种顺序小文件足够，遍历目录时superblock（块 1）和 inode 区会反复换入换出，每读一个 dirent 的inode 都可能多一次 SD 读。xv6 的 buffer cache 是 30 块 LRU，这里刻意省掉了。

读路径（无 dcache、无 icache，每次 lookup 全程走磁盘）：

\begin{txtcode}
namei("/bin/sh")
+-- 逐分量 directory_lookup()
    +-- read_inode(dir_inum)                 块 = inodestart + inum/16
    +-- inode_read() 逐 dirent              inode_block(): 前 12 块直接, 之后查间接块
    +-- names_equal(disk_name, component)   名字 ≤30 字节, NUL/长度双判
\end{txtcode}

\section{bootfs：\file{/boot} 的 FAT32 交接}

\begin{note}
现在只是后备路径：正常情况下 \file{/boot} 由 \code{fat32.c} 挂载，见第~\ref{chap:bootfat32} 章。
\end{note}

背景：XR829 三个固件文件现在放在 FAT32 分区里（方便换固件不用重烧bootpack），但内核的 minifs 是只读小文件系统、rootdev 只认分区 2。解法是分工：\textbf{SPL 读 FAT32（fat32.c，MBR + BPB + 簇链，支持长文件名），内核只认 DDR 里已经就位的镜像}。

SPL 把固件装到 DDR 固定地址后，向 bootpack 头补写三个条目（\code{offset=0}表示"常驻镜像，非 pack 内偏移"），重算头 CRC：

\begin{table}[H]
\centering
\small
\begin{tabular}{llll}
\toprule
文件 & bootpack 条目 & DDR 地址 & 容量上限 \\
\midrule
boot\_xr829.bin & \code{boot\_xr829} & 0x40100000 & 64 KiB \\
fw\_xr829.bin & \code{fw\_xr829} & 0x40110000 & 192 KiB \\
sdd\_xr829.bin & \code{sdd\_xr829} & 0x40140000 & 64 KiB \\
\bottomrule
\end{tabular}
\end{table}

内核侧 \code{bootfs.c} 把这三条暴露成普通文件：

\begin{itemize}
  \item \code{lookup("/boot/fw\_xr829.bin")} → \code{t113\_bootpack\_find("fw\_xr829")}，\code{node.id = load\_addr}（物理地址当文件 id）；
  \item \code{read()} 就是 \code{memcpy}——数据已在 DDR，零拷贝语义。
\end{itemize}

XR829 驱动的 \code{need\_image()} 因此从硬编码 bootpack 查询改成了标准VFS 调用（\code{t113\_vfs\_open("/boot/fw\_xr829.bin")} → 读到页分配的缓冲），驱动不再依赖 bootloader 私有的地址常量。若 FAT32 分区缺失，SPL 回退使用 bootpack 内打包的固件（老路径仍保留）。

\section{devfs：\file{/dev} 设备文件}

\code{src/kernel/fs/devfs.c} 对应 Linux 早期的 devfs（现代等价物是devtmpfs）：

\begin{itemize}
  \item \code{t113\_devfs\_init()}：挂到 \file{/dev}，16 个节点槽，设备名 ≤ 30 字节且不含 \file{/}；
  \item 驱动 probe 时做两件事：\code{t113\_chrdev\_register(major, minor, ops)} 把设备号和驱动操作登记进 \code{chrdev.c} 的 \code{chrdev\_table[32]}（\code{ops} 至少要有read 或 write，设备号去重）；\code{t113\_devfs\_create\_chrdev(name, major, minor)}在 \file{/dev} 下建名字（名字去重）。这和 Linux 的 \code{register\_chrdev\_region} + \code{cdev\_add}、再由 devtmpfs 建节点是同一个分工；
  \item \code{lookup("/dev/console")}：\code{node.id = (major<<16)|minor}、\code{type = DEVFS\_T\_CHAR}、\code{private\_data} 指向节点；\file{/dev} 本身作为目录节点（type 1）返回，\code{readdir("/dev", i)} 枚举已注册设备；
  \item \code{read}/\code{write} 按 \code{node.id} 调 \code{t113\_chrdev\_read/write()}，再由\code{chrdev\_table} 找到驱动的 \code{ops}，\textbf{忽略 offset}——字符设备没有seek 语义，这和 Linux \code{struct file\_operations} 里 \code{.llseek = no\_llseek} 的约定一致。
\end{itemize}

当前唯一注册者是 UART 平台驱动（\code{uart.c} probe）：\code{console}（major 1, minor 0）和 \code{ttyS0}（major 3, minor 0），major 取值沿 Linux 惯例（1 = 内存/终端类，3 = 第一个 8250 串口）。注册时序有要求：\code{kernel\_boot()} 里的 \code{t113\_devfs\_init()} 必须早于驱动绑定 \code{t113\_driver\_core\_init()}，否则 probe 里的注册会失败。

\textbf{用户态 stdio}：\file{/init} 第一件事是 \code{open("/dev/console", O\_RDWR)} 得到fd 0，再 \code{dup(0)} 两次得到 fd 1、2；\code{syscall.c} 不再特判这三个 fd。shell继承的 stdin/stdout 就是这个 devfs 文件，三个 fd 共享同一个 \code{struct file}。另一个注册者是 \code{drivers/input/gpio\_keys.c} 的 \file{/dev/key0}。

\section{用户态视角：从 \file{/init} 到 shell}

系统调用到 VFS 的映射（\code{src/kernel/syscall.c}，每进程\code{files[T113\_MAX\_FDS=16]}，fd 0/1/2 由 \file{/init} 打开 \file{/dev/console} 建立）：

\begin{table}[H]
\centering
\small
\begin{tabularx}{\linewidth}{>{\raggedright\arraybackslash}X>{\raggedright\arraybackslash}X>{\raggedright\arraybackslash}X}
\toprule
系统调用 & 内核路径 & 备注 \\
\midrule
\code{open(path, flags)} & \code{sys\_open} → \code{t113\_vfs\_open} & 用户指针逐字节经 \code{user\_vm\_range\_mapped} 校验拷入 \\
\code{read(fd, buf, n)} & \code{t113\_vfs\_read} & 控制台经 devfs → chrdev → UART \\
\code{write(fd, buf, n)} & \code{t113\_vfs\_write} & 单次 ≤ 256 B \\
\code{dup(fd)} & \code{t113\_vfs\_file\_get} & 新 fd 共享同一个 \code{struct file} \\
\code{close(fd)} & \code{t113\_vfs\_close} & 引用计数归零才释放槽位 \\
\code{readdir(path, i)} & \code{t113\_vfs\_readdir} & \code{ls} 的实现基础 \\
\code{exec(path, argv)} & \code{sys\_exec} → \code{t113\_vfs\_open} + ELF 加载 & 镜像 ≤ 16 页（64 KiB） \\
\code{getcwd} & 常量 \code{"/"} & 尚无每进程 cwd，pwd 是假的 \\
\bottomrule
\end{tabularx}
\end{table}

内核加载 \file{/init} 本身（\code{kernel\_main} 尾部）就是一次普通 VFS 读：open → 按文件大小 \code{alloc\_pages} → read → \code{create\_init\_user\_process} → 释放临时页。\file{/init} 再 \code{fork()} + \code{exec("/bin/sh")}，shell 的每次命令执行又是一次 \code{exec} → VFS → minifs 读块。整条用户态生命线都吊在 rootdev 的 CMD17 上。

\section{与 Linux / xv6 的对照}

\begin{table}[H]
\centering
\small
\begin{tabularx}{\linewidth}{>{\raggedright\arraybackslash}X>{\raggedright\arraybackslash}X>{\raggedright\arraybackslash}X}
\toprule
miniOS & xv6 & Linux \\
\midrule
\code{vfs\_mount} 树 + 前缀匹配 & 无（只有一个 fs） & 挂载树 + namespace \\
\code{super\_block} / \code{inode} / \code{dentry} 对象池 & \code{itable} inode 缓存，无 dentry & super\_block、icache、dcache \\
\code{inode\_operations} / \code{file\_operations}（经 \code{vfs\_backend\_ops} 适配） & 无，直接调 fs 函数 & \code{inode\_operations} + \code{file\_operations} + \code{super\_operations} \\
\code{rootdev}（MBR 分区） & \code{ide} 驱动 & 块设备层 + 分区扫描（\code{read\_partition}） \\
minifs 只读 & xv6 fs（含 log 写入） & ext4 等 \\
单块缓存 & 30 块 buffer cache & page cache \\
\file{/boot} FAT32（失败时退回 bootfs） & — & \file{/boot} 分区 + firmware loader \\
devfs + \code{chrdev\_register} & mknod + \code{devsw[]} & devtmpfs + \code{device\_create}（+ udev 用户态改名） \\
\code{kernel\_main} 读 \file{/init} & 第一个进程硬编码在 initcode.S & initramfs 的 \file{/init} → exec \file{/sbin/init} \\
PANIC + \code{wfi} & \code{panic()} & \code{panic()} \\
\bottomrule
\end{tabularx}
\end{table}

\section{构建与烧写}

\begin{shcode}
make t113                    # 内核 + minifs.img（mkfs_minifs 打包 init/sh/ls/...）
make install T113_DISK=/dev/diskX
# 内部依次：
#   dd minios-t113.bin  @ 8 KiB    (Boot0)
#   dd bootpack.bin     @ 128 KiB  (kernel)
#   cp firmware/xr829/*.bin → /Volumes/bootfs   (FAT32, install-firmware)
#   dd minifs.img → /dev/diskXs2                (install-rootfs)
\end{shcode}

换固件只动 FAT32 分区；改用户程序（\code{user-t113/*.c}）只需重跑\code{install-rootfs}，内核和 bootpack 不用重烧。

\section{预期串口输出}

（LBA/大小以实际卡为准，下面是格式示例）

\begin{txtcode}[basicstyle=\ttfamily\scriptsize]
SPL: SDHC card ready
SPL: load kernel -> 0x40200000 bytes=0x?????? ok
SPL: FAT32 partition LBA=0x00002000
SPL: FAT32 load boot_xr829.bin bytes=0x888
SPL: FAT32 load fw_xr829.bin bytes=0x2aa20
SPL: FAT32 load sdd_xr829.bin bytes=0x2e8
...
bootpack: kernel @0x40200000 bytes=0x????????
bootpack: boot_xr829 @0x40100000 bytes=0x00000888
bootpack: fw_xr829 @0x40110000 bytes=0x0002aa20
bootpack: sdd_xr829 @0x40140000 bytes=0x000002e8
rootdev: MBR partition 2 type=0x83 lba=0x?????? sectors=0x??????
minifs: mounted block rootfs blocks=0x8000 inodes=0xc8
fat32: /boot mounted rw, partition lba=0x???????? clusters=0x???????? cluster-bytes=0x????????
...
init: starting shell on console
\end{txtcode}

\code{rootdev:}、\code{minifs:} 先出现，然后才是 \code{fat32:}：\code{kernel\_boot()} 先挂 \file{/}（它顺带初始化 SD 卡），再挂 \file{/boot}。FAT32 挂载失败时会多一行 \code{fat32: mount failed, using SPL handoff images}，接着是 \code{bootfs: FAT32 handoff mounted at /boot}。devfs 挂载没有打印，驱动绑定后 \code{platform\_dev\_info} 会打印 uart 绑定信息。
